% !TeX root = ../main.tex
\chapter{Conclusions and Future Work}
\label{ch:conclusions}
\section{Purpose and Completed Contribution}
This project investigates how a digital-twin-oriented training environment can
replace manual driving demonstrations when developing a navigation policy for
a resource-constrained board. The implemented solution assigns different roles
to different models: a PPO-trained neural teacher learns from reward feedback
in a virtual vehicle, while a supervised tree student learns the teacher's
LiDAR-to-action mapping. Training remains a host-side activity and the exported
student is the candidate for eventual onboard inference.

The contribution is the connection between these stages. The virtual platform
provides controlled interaction and privileged state; the teacher supplies
commands without human per-sample action labels; the collector preserves scan
and action pairing; and the student/export workflow produces a separately
evaluable controller. The present evidence covers the virtual training side,
rather than a completed physical twin or a demonstrated low-power deployment.

\section{Answers to the Research Questions}
\textbf{RQ1: Automated supervision.}
The selected PPO--MLP teacher supplies two completed collection batches containing
1,439,983 and 1,079,979 original LiDAR--action pairs, respectively. The collection
process uses teacher commands rather than human driving labels, with model,
map and episode provenance retained. Its separate T3 qualification completes
eight episodes without recorded collision or listed motion failure. These
records establish automated supervision under the implemented conditions.
They do not measure how much operator time is saved or show that teacher labels
are better than human labels. Reward design, modelling and experiment selection
remain engineering responsibilities.

\textbf{RQ2: Transfer to compact policies.}
The automatically generated data support useful student cruising. Across the
normal and scattered-dropout conditions, the forests complete between 57 and
59 of 60 primary episodes and the single tree completes all sixty. Over all
conditions, the forests complete 62--67 of 80 episodes and the single tree
completes 74. Offline action accuracy does not reproduce this closed-loop
ranking. Contiguous sector loss accounts for 46 of the 53 primary collisions,
and the selected scattered-dropout augmentation gives no consistent advantage.
The strongest tested candidate is therefore the single-tree configuration,
although the comparison does not isolate one structural cause. This result
supports the wider objective of obtaining a compact useful policy, while
qualifying the initial choice of a random forest.

\textbf{RQ3: Resource-constrained deployment.}
Tree depth and estimator count make the threshold-evaluation workload explicit,
and the exported C++ implementation agrees with the fitted model on the checked
host inputs. These properties justify pursuing embedded integration. They do
not establish Flash occupancy, peak RAM, worst-case control latency or energy
per decision on the intended board. Nor does successful export establish that
the physical sensor and actuator interfaces reproduce the simulated ones.
The final low-power deployment objective therefore remains partially answered.

\section{Interpretation of the Method Choices}
RL addresses the lack of an initial action-labelled dataset by learning from
sequential outcomes. PPO provides the implemented actor--critic training
procedure, and the MLP represents the teacher's policy and value function.
Privileged geometry is useful because the teacher's immediate role is to
produce demonstrations; the student is subsequently restricted to the sensor
representation available for deployment. A forest offers bounded-depth tree
inference and ensemble aggregation, while the single tree tests whether a
simpler configuration suffices. These choices form a coherent division of
training and deployment responsibilities, not a claim that PPO, MLPs or forests
are individually optimal for every navigation task.

The experimental lesson is that the workflow should be judged by the quality
of the generated supervision and the behaviour and implementation cost of its
student. Selecting a family for practical reasons is a starting hypothesis;
the closed-loop comparison can legitimately favour a different compact member
of the candidate set. The main research purpose is preserved when that happens.

\section{Limitations and Priorities}
The next simulation priority is the frozen-candidate comparison on independent
maps and perturbed student starts. Recovery criteria and observation windows
should be fixed before execution. Selected failure replays with original and
corrupted scans, mask indices and decision timestamps would then help explain
the recurrent R1 and sector collisions. Matched structural comparisons would
be needed before attributing the tree/forest difference to depth or ensembling.

A direct test of the manual-demonstration replacement claim would compare the
automated workflow with a documented human-driven or scripted-teacher baseline
under matched tasks and student settings. Relevant measurements include
operator involvement, collection cost, coverage and resulting driving behaviour.
The absence of that comparison limits claims about efficiency, without changing
the recorded fact that this pipeline generated its action labels automatically.

The later physical phase should calibrate the virtual vehicle against the real
platform, verify sensor/action correspondence, and measure compiled memory,
control latency and energy together with repeated driving outcomes. This would
connect the implemented virtual training side to the intended physical twin and
low-power controller. Until then, the dissertation establishes an automated
simulation-to-student training mechanism and identifies the evidence still
needed for its full embedded objective.
